% ============================================================================
% sec5_empirical_strategy.tex — REVISION-V6: Empirical strategy (RAND reframe)
% Paper: Frequent Losers as Cover Bidders in Public Procurement
% Target: RAND Journal of Economics
% ============================================================================

\section{Empirical Strategy}
\label{sec:empirical}

Our empirical strategy has three tiers, ordered by the strength of causal
claims. \emph{First}, we estimate the structural model of cover-bidder
deployment from Section~\ref{sec:structural}, which identifies the
cover-bid distribution and implied markup directly from the bid-level
data (Section~\ref{sec:structural_estimation}). \emph{Second}, we
provide reduced-form evidence using high-dimensional OLS with item, year,
and purchasing-unit fixed effects, supplemented by a leave-one-out
instrumental variable (Section~\ref{sec:reduced_form}). \emph{Third}, we
validate the FL screen's detection performance against CADE cartel
convictions using ROC analysis (Section~\ref{sec:detection_validation}).
Non-structural corroboration is provided by corrected Bajari--Ye
bid-coordination tests (Section~\ref{sec:bajari_ye}).

% ── 5.1 Structural Estimation ────────────────────────────────────

\subsection{Structural Likelihood Estimation}
\label{sec:structural_estimation}

We estimate the mixture model specified in % R1-FIX: sample-size pipeline clarified
Section~\ref{sec:likelihood} using the bid-level data. The full
dataset contains approximately 40 million bids from 41,000 firms.
Of these, 39.7 million are non-FL bids and 0.2 million are FL bids.
After restricting to \emph{losing} bids with valid prices and
item-year fixed-effect coverage, the estimation sample comprises
23.2 million genuine (non-FL) losing bids and 189,381 FL losing bids
(Table~\ref{tab:structural_params}). The estimation proceeds in two
stages.

\paragraph{Stage 1: Genuine bid distribution.}
We estimate the genuine bid parameters $(\boldsymbol{\beta},
\sigma_g^2)$ from the 23.2 million non-FL losing bids using maximum
likelihood under the log-normal specification~\eqref{eq:genuine_dist}:
\begin{equation}
  (\hat{\boldsymbol{\beta}}, \hat{\sigma}_g^2) =
  \operatorname*{arg\,max}_{\boldsymbol{\beta}, \sigma_g^2}
  \sum_{t=1}^{T} \sum_{j \in \mathcal{G}_t}
  \log \phi\!\left(
    \frac{\log b_{jt} - \mathbf{x}_j'\boldsymbol{\beta}
    - \hat{\alpha}_i - \hat{\lambda}_t}{\sigma_g}
  \right)
  \label{eq:mle_genuine}
\end{equation}
where $\phi(\cdot)$ is the standard normal density and
$\hat{\alpha}_i$, $\hat{\lambda}_t$ are absorbed as fixed effects.
In practice, this reduces to an OLS regression of $\log b_{jt}$ on
$\mathbf{x}_j$ with item and year fixed effects, from which we obtain
$\hat{\sigma}_g^2$ as the residual variance.

\paragraph{Stage 2: Cover bid distribution.}
Conditional on the winning bid $b^*_t$ in each tender, we estimate
the cover-bid parameters from the 2.5 million FL bids. Under
Regime~1:
\begin{equation}
  \hat{\delta} = \max_{j \in \mathcal{F}_t, t}
  (b_{jt} - b^*_t) \quad \text{(maximum spread)}
  \label{eq:mle_delta}
\end{equation}
Under Regime~2:
\begin{equation}
  \hat{\sigma}_c^2 = \frac{1}{|\mathcal{F}|}
  \sum_{t=1}^{T} \sum_{j \in \mathcal{F}_t}
  \left(\log b_{jt} - \log b^*_t - \hat{\epsilon}\right)^2
  \label{eq:mle_sigma_c}
\end{equation}
where $\hat{\epsilon} = \overline{\log b_{jt} - \log b^*_t}$ is the
mean log-spread of FL bids above the winning bid.

\paragraph{Model selection.}
We select between regimes using the Bayesian Information Criterion
(BIC) and a Kolmogorov--Smirnov goodness-of-fit test:
\begin{equation}
  \text{BIC}_r = -2\ell_r + k_r \log N, \quad r \in \{1, 2\}
  \label{eq:bic}
\end{equation}
where $\ell_r$ is the maximized log-likelihood under regime $r$,
$k_r$ is the number of parameters, and $N$ is the number of FL bids.
The regime with lower BIC is selected. Additionally, we compute the KS
distance between the fitted cover-bid CDF and the empirical FL bid
distribution; the regime with smaller KS distance provides a better
fit to the data.

\paragraph{Implied markup.}
The structural markup is computed as:
\begin{equation}
  \hat{\mu}_{\text{structural}} =
  \frac{E[b^*_t \mid \text{FL present}] -
        E[b^*_t \mid \text{FL absent}]}{E[b^*_t \mid \text{FL absent}]}
  \label{eq:structural_markup}
\end{equation}
where the expectations are taken over the estimated model. This provides
a structural counterpart to the reduced-form OLS coefficient
($\hat{\beta} = 0.064$) and the IV estimate ($\hat{\beta}_{\text{IV}} =
0.194$). If the structural model is correctly specified, the implied
markup should fall within the OLS--IV range.

\paragraph{Standard errors.}
We compute standard errors via parametric bootstrap with 500
replications, resampling tenders (not individual bids) to preserve the
within-tender correlation structure. Each bootstrap iteration
re-estimates both stages of the model. With 16 CPU cores, the bootstrap
completes in approximately 4--8 hours. % REVISION-V6: 500 replications

% ── 5.2 Reduced-Form Identification ──────────────────────────────

\subsection{Reduced-Form Identification}
\label{sec:reduced_form}

\paragraph{OLS baseline.}
Our baseline specification estimates the conditional association between
FL presence and tender outcomes:
\begin{equation}
  y_{igt} = \beta \cdot \text{losers}_{igt} + \mathbf{x}_{igt}'
  \boldsymbol{\delta} + \alpha_g + \lambda_t + \gamma_k + \varepsilon_{igt}
  \label{eq:ols_main}
\end{equation}
where $y_{igt}$ is the outcome for tender-item $i$ in item group $g$ at
time $t$ and purchasing unit $k$; $\text{losers}_{igt} \in \{0,1\}$
indicates FL presence; $\mathbf{x}_{igt}$ includes a convite indicator;
$\alpha_g$, $\lambda_t$, and $\gamma_k$ are item, year, and PBU fixed
effects; and $\varepsilon_{igt}$ is the error term clustered at the item
level.

The OLS coefficient $\hat{\beta}$ identifies a conditional
association---the average difference in outcomes between FL-present and
FL-absent tenders within the same item type, year, and purchasing unit.
Under Prediction~\ref{pred:price_v6}, $\hat{\beta} > 0$ for prices. The
OLS estimate of 0.064 (6.4\%) represents a lower bound on the causal
effect if the FL indicator contains measurement error (binary treatment
misclassification). We assess the sensitivity of $\hat{\beta}$ to
unobserved confounders using the \citet{cinelli2020making} framework,
which yields a robustness value of $RV_{q=1} = 17.5\%$: an unobserved
confounder would need to explain at least 17.5\% of the residual
variation in both FL presence and log prices to reduce $\hat{\beta}$ to
zero.

\paragraph{Instrumental variable.}
To bracket the OLS estimate from above, we instrument FL presence with % REVISION-V6: IV reframed as upper bound, not primary
a leave-one-out measure of FL supply:
\begin{equation}
  Z_{kgt} = \sum_{j \neq k} \mathbf{1}[\text{FL firm active at PBU } j
  \text{ in group } g, \text{ year } t]
  \label{eq:iv_v6}
\end{equation}
The identifying assumption is that FL activity at other PBUs in the same
product market and year affects focal-PBU prices only through the
probability of FL participation---a supply-side exclusion restriction in
the spirit of \citet{goldsmithpinkham2020bartik}. We treat the IV
estimate as a Local Average Treatment Effect for complier tenders: an
upper bound on the average price effect that corrects for
measurement-error attenuation in the FL indicator. The IV is not our
primary identification strategy; we report it to bracket the OLS
estimate and verify that the structural model's implied markup falls
within the OLS--IV range.

\paragraph{Balance and specification sensitivity.}
Balance tests (Table~\ref{tab:iv_balance} in
Appendix~\ref{app:robustness}) reveal statistically significant
differences across all four observable characteristics, a pattern that
is more concerning than any single coefficient. While the economic
magnitudes are modest (all standardized differences below
$0.03\sigma$), we do not dismiss this imbalance. The sensitivity
analysis and matching estimates
(Section~\ref{sec:robustness}) bound the extent to which unobserved
selection could explain the price effect. We report the IV for
completeness and as a bracketing device, not as the basis for our
causal claims.

% ── 5.3 Detection Validation ─────────────────────────────────────

\subsection{Detection Performance Validation}
\label{sec:detection_validation}

We evaluate the FL screen as a detection tool using CADE cartel
convictions as ground truth. The validation uses the temporal split
described in Section~\ref{sec:cade_sample}: FL status is defined using
2009--2014 data and evaluated against 2015--2019 CADE outcomes.

\paragraph{ROC analysis.}
We vary the IQR multiplier from 0.5 to 5.0 in steps of 0.1, producing
46 different FL thresholds. For each threshold, we compute:
\begin{align}
  \text{TPR}(m) &= \frac{|\{j : \text{FL}_m(j)=1\} \cap
    \{j : \text{CADE}(j)=1\}|}{|\{j : \text{CADE}(j)=1\}|}
  \label{eq:tpr} \\
  \text{FPR}(m) &= \frac{|\{j : \text{FL}_m(j)=1\} \cap
    \{j : \text{CADE}(j)=0\}|}{|\{j : \text{CADE}(j)=0\}|}
  \label{eq:fpr}
\end{align}
where $\text{FL}_m(j)$ indicates whether firm $j$ is classified as FL
under multiplier $m$, and $\text{CADE}(j)$ indicates whether firm $j$
co-participates with a CADE-convicted cartelists. The area under the ROC
curve (AUC) measures the FL screen's overall discrimination ability; AUC
= 0.5 corresponds to random classification and AUC = 1.0 to perfect
discrimination.

\paragraph{Optimal threshold.}
We determine the welfare-maximizing IQR multiplier using Youden's $J$
statistic:
\begin{equation}
  m^*_{\text{optimal}} = \operatorname*{arg\,max}_m \;
  \left[\text{TPR}(m) - \text{FPR}(m)\right]
  \label{eq:youden}
\end{equation}
which balances true-positive and false-positive rates. The optimal
threshold provides a data-driven calibration of the statistical cutoff
introduced in Section~\ref{sec:fl_definition}, connecting the ad-hoc
IQR rule to a formal detection-theoretic criterion.

\paragraph{Comparison with Imhof screens.}
We compare the FL screen with a proxy for the \citet{imhof2019detecting}
bid-level screen, implemented as a within-item-group coefficient of
variation (CV) median split. For each feature set---FL indicators only,
Imhof-style bid-level features only, and the combination---we train a
random forest classifier using 5-fold stratified cross-validation and
report the mean AUC. The incremental AUC of combining FL with bid-level
features measures the complementarity of the two screening approaches.

% ── 5.4 Bajari-Ye Tests ──────────────────────────────────────────

\subsection{Bajari--Ye Bid Coordination Tests}
\label{sec:bajari_ye}

We implement the \citet{bajari2003deciding} tests for exchangeability
and conditional independence using corrected first-stage residuals.

\paragraph{First stage (corrected).} % REVISION-V6: n_bids excluded, noted explicitly
The first stage regresses $\log$ bid values on observable cost shifters
only:
\begin{equation}
  \log b_{jt} = \mathbf{x}_j'\boldsymbol{\phi} + \alpha_i + \lambda_t
  + u_{jt}
  \label{eq:by_first_stage}
\end{equation}
where $\mathbf{x}_j$ includes firm size (porte), firm age, and CNAE
sector dummies. Critically, we exclude the number of bids in the
tender ($n_{\text{bids}}$) from the first stage. As
\citet{bajari2003deciding} note (p.~978), the first-stage regressors
must be cost shifters, not equilibrium market outcomes. The number of
bids is jointly determined with bid levels and is mechanically increased
by FL presence (Table~\ref{tab:nfirms_excl}); including it would
generate spurious correlation between the residuals and FL status.

\paragraph{Exchangeability test.}
Using the first-stage residuals $\hat{u}_{jt}$, we conduct a
Kolmogorov--Smirnov test comparing the distributions of FL and non-FL
residuals. Under competitive bidding with symmetric information, all
bidders' residuals should come from the same distribution
(Prediction~\ref{pred:exchange_v6}).

\paragraph{Conditional independence test.}
For each tender with $\geq 2$ FL bids, we compute the mean product of
pairwise FL residuals:
$\bar{r}_t = \binom{m_t}{2}^{-1} \sum_{j < k}
\hat{u}_{jt} \hat{u}_{kt}$.
Under conditional independence, $E[\bar{r}_t] = 0$. We compare the FL
pairwise product with the non-FL pairwise product using a tender-level
bootstrap (1,000 iterations).

\paragraph{Reframing the fake-groups placebo.}
The fake-groups placebo randomly splits non-FL firms into two groups
within each tender and computes within-group pairwise products. This
test does \emph{not} establish a zero null: common tender-level shocks
(reference price, commodity prices, PBU-specific factors) induce
positive residual correlations among \emph{all} bidders in the same
tender, regardless of FL status. The economically meaningful comparison
is the FL-vs-non-FL pairwise \emph{difference} (bootstrap 95\% CI
excludes zero), which isolates the within-tender excess correlation of
FL bids above and beyond common shocks.

\paragraph{Tender fixed effects.}
As a more demanding specification, we re-estimate the first-stage bid
equation with tender $\times$ item fixed effects, which absorb all
tender-level shocks. Under tender FE, the FL pairwise product drops
relative to the baseline specification. We report both variants and
interpret the tender-FE results as the most conservative test of
within-tender FL bid coordination.

% ── 5.5 Network-Split Heterogeneity ──────────────────────────────

\subsection{Network-Split Heterogeneity Test} % R1-FIX: added missing subsection (referee Major #4)
\label{sec:empirical_network}

To test Proposition~\ref{prop:market_selection}---that cover bidding
concentrates in competitive markets---we split FL firms into two
subgroups based on co-bidding network characteristics.

\paragraph{Winner HHI.}
For each FL firm $j$, we compute the Herfindahl--Hirschman Index of
winning firms across all tenders in which $j$ participates:
\begin{equation}
  \text{HHI}_j = \sum_{w=1}^{W_j} s_{wj}^2
  \label{eq:winner_hhi}
\end{equation}
where $s_{wj}$ is the share of firm $j$'s tenders won by winner $w$,
and $W_j$ is the number of distinct winners. A high $\text{HHI}_j$
indicates that $j$ operates in markets dominated by a single or few
winners (concentrated markets); a low $\text{HHI}_j$ indicates
diversified winner patterns (competitive markets).

\paragraph{Classification.}
We classify FL firms as \emph{concentrated-market} if
$\text{HHI}_j$ exceeds the sample median, and
\emph{competitive-market} otherwise. The median split avoids ad hoc
threshold selection while maintaining approximate sample balance.
We then estimate the baseline price regression
(Equation~\ref{eq:ols_main}) separately for each subgroup,
interacting $\text{losers}_{igt}$ with the market-structure
indicator.

\paragraph{Interpretation.}
Under Proposition~\ref{prop:market_selection}, the FL--price effect
should be concentrated in competitive-market FL firms (low HHI), where
the marginal return to manufactured competition is highest, and absent
in concentrated-market FL firms, where the cartel can sustain prices
through market power alone.

% REVISION-V6: RDD removed. Aggregating unit prices × quantities to
% contract-level totals reveals that convite assignment does not exhibit
% a discontinuity at R$80K.
